%% ma: L12-B15-0024
%% lop: 12
%% chuong: 5
%% bai: 15
%% dang: 12.15.19
%% loai: TN4
%% muc: TH
%% dap_an: "D"
%% nguon: "(NBV)-10 ĐỀ KIỂM TRA HK2-TOÁN 12-DUNG CHUNG 3 SÁCH.pdf (H04, Đề số 4), Phần 1 Câu 2"
%% kien_thuc: "Bài 15 (phương trình tham số của đường thẳng đi qua hai điểm)"
%% trang_thai: cho-duyet
%% ly_do: "Dạng mới đề xuất 12.15.19 (Phương trình tham số của đường thẳng đi qua hai điểm), chờ giáo viên duyệt."
%% ghi_chu: "Đáp án gốc D đúng."
\begin{ex}
Trong không gian $Oxyz$, cho hai điểm $M(1;0;1)$ và $N(3;2;-1)$. Đường thẳng $MN$ có phương trình tham số là
\choice{$\begin{cases}x=1+2t\\y=2t\\z=1+t\end{cases}$}{$\begin{cases}x=1+t\\y=t\\z=1+t\end{cases}$}{$\begin{cases}x=1-t\\y=t\\z=1+t\end{cases}$}{\True $\begin{cases}x=1+t\\y=t\\z=1-t\end{cases}$}
\loigiai{$\overrightarrow{MN}=(2;2;-2)=2(1;1;-1)$. Đường thẳng đi qua $M(1;0;1)$ với vectơ chỉ phương $(1;1;-1)$ có phương trình $x=1+t$, $y=t$, $z=1-t$. Chọn D.}
\end{ex}
